\section{Limitations}
The oracle evaluation measures what the manifest permits along each task's ground-truth tool sequence, not what a model
does. A model may solve a task with other tools, or follow an injection only partly, so \Q1--\Q3 bound attack success
rather than measure it; \Q4 measures it. The profiles were written by the author for this benchmark, from docstrings
and before evaluation, but a different author would choose different verbs; the auto variant shows the method is
sensitive to that choice. The matcher is lexical: a paraphrase without a listed verb (``get Bob his money back'') is
missed, and an instruction that names an action it forbids (``do not delete anything'') opens it. Knowing the
profiles does not let an attacker widen the manifest, because it depends only on the user's text, but it does tell
them which tools stay open. A user task that delegates its instructions to content is unprotected by design. Results are for AgentDojo's four simulated suites and
English instructions only.

\section{Conclusion}
A tool whitelist computed once from the user's own words, with no model call, blocks most prompt-injection attacks
that need a capability the user did not ask for, and keeps most legitimate tasks working. Its failures are
informative. Benign tasks break when the user delegates the instruction to untrusted content, and \N{unblockable}\% of
attacks cannot be stopped by any tool-level policy because they reuse the capabilities the user's task requires. For
cost-sensitive deployments the manifest is a sensible first layer. The cases it cannot handle, delegation and reuse of
the user's own capabilities, are exactly where model-based or data-flow defenses should spend their cost.
